\introduction

Актуальность исследования обусловлена широким распространением облачных технологий, контейнеризации и микросервисной архитектуры в современных информационных системах. Переход от монолитных приложений к распределённым сервисам позволил повысить масштабируемость и отказоустойчивость программных решений, однако одновременно привёл к усложнению процессов администрирования и обеспечения безопасности \cite{Dragoni_2017_Microservices_yesterday_today_and_tomorrow}.

В настоящее время контейнеризация и оркестрация контейнеров являются базовыми технологиями при построении корпоративных и облачных систем. Одним из наиболее распространённых инструментов управления контейнеризированными приложениями выступает платформа Kubernetes, обеспечивающая автоматическое развертывание, масштабирование и координацию контейнеров в распределённой инфраструктуре. Широкое применение Kubernetes в промышленной эксплуатации подтверждает его значимость для современных информационных систем \cite{Container_Security_Extensive_Roadmap}. Вместе с тем динамическая структура Kubernetes-кластера создаёт дополнительные риски в области информационной безопасности.

Одним из наиболее информативных источников данных для анализа безопасности Kubernetes-кластера являются audit-логи, формируемые компонентом API-сервера. Указанные журналы содержат сведения о запросах к объектам кластера, включая тип операции, пользователя, целевой ресурс, код ответа и иные параметры запроса. Анализ таких данных позволяет выявлять нетипичные сценарии поведения, которые могут указывать на наличие атак, нарушений политики безопасности или ошибок конфигурации \cite{KubernetesAudit}.

Однако использование audit-логов затрудняется высокой интенсивностью их формирования, значительным объёмом данных и сложностью выявления скрытых зависимостей между событиями. Традиционные сигнатурные и правил-ориентированные методы обнаружения аномалий недостаточно эффективны в условиях изменяющегося характера атак и разнообразия легитимного поведения пользователей. Кроме того, audit-логи представляют собой последовательности событий, в которых важную роль играет не только факт отдельного обращения, но и порядок, контекст и временные зависимости между событиями. В связи с этим актуальным является применение методов машинного обучения и нейронных сетей, способных автоматически выявлять аномальные последовательности и скрытые закономерности в данных.

Таким образом, актуальность настоящей работы определяется следующими факторами:

\begin{enumerate}
    \item ростом числа атак на облачные и контейнерные инфраструктуры, включая Kubernetes-кластеры \cite{Kubernetes_Security_Report};
    \item высокой интенсивностью формирования audit-логов и сложностью их ручного анализа в распределённых системах;
    \item недостаточной эффективностью традиционных сигнатурных и правил-ориентированных методов при выявлении новых и нетипичных сценариев атак;
    \item необходимостью разработки автоматизированных средств обнаружения аномалий, учитывающих временную структуру и последовательный характер audit-событий в Kubernetes-кластере.
\end{enumerate}

Целью выпускной квалификационной работы является разработка средства обнаружения аномалий в Kubernetes-кластере на основе методов нейронных сетей, обеспечивающего высокую точность выявления аномальных событий.

Для достижения поставленной цели необходимо решить следующие задачи:

\begin{enumerate}
    \item провести обзор и анализ научной литературы в области обнаружения аномалий в контейнерных средах;  
    \item рассмотреть и сравнить существующие подходы и методы обнаружения аномалий в контейнерных средах на основе нейронных сетей;
    \item выполнить моделирование одного из подходов обнаружения аномалий;
    \item разработать архитектуру средства обнаружения аномалий в Kubernetes-кластере;
    \item сформулировать математическую постановку задачи 
    \item сформировать или получить набор данных для обучения модели нейронной сети;
    \item построить и обучить модель нейронной сети;
    \item разработать программное обеспечение на основе выбранной модели НС;
    \item тестирование разработанного средства;
    \item обоснование экономической целесообразности разработанного средства;
    \item правовое обоснование выполнения выпускной квалификационной работы. 
\end{enumerate}
